\chapter{Malaria Testing}

\section*{What Is This Test?}
Malaria testing detects \textit{Plasmodium} species parasites, the causative agents of malaria, a life-threatening mosquito-borne illness transmitted by female \textit{Anopheles} mosquitoes. Five species infect humans: \textit{Plasmodium falciparum} (most virulent, responsible for >90\% of malaria deaths, especially in sub-Saharan Africa), \textit{P. vivax} (most geographically widespread, with dormant hypnozoite liver stages causing relapses), \textit{P. ovale} (two subspecies: \textit{ovale curtisi} and \textit{ovale wallikeri}, also with hypnozoites), \textit{P. malariae} (causes chronic, low-level parasitemia lasting decades), and \textit{P. knowlesi} (a zoonosis from macaque monkeys in Southeast Asia, morphologically similar to \textit{P. malariae} on microscopy). Malaria is a medical emergency in non-immune travelers. \textit{P. falciparum} can progress to severe malaria (cerebral malaria, severe anemia, acute kidney injury, ARDS, DIC, metabolic acidosis, hypoglycemia, and multi-organ failure) within 24--48 hours, with mortality of 15--20\% despite ICU care. Uncomplicated malaria presents with paroxysmal fever/chills/sweats (every 48 hours for \textit{P. vivax} and \textit{P. ovale}/tertian, every 72 hours for \textit{P. malariae}/quartan, irregular for \textit{P. falciparum}), headache, myalgias, nausea, hepatosplenomegaly, anemia, and thrombocytopenia \cite{white2014malaria}. Diagnosis relies on microscopic examination of Giemsa-stained thick and thin blood films and rapid diagnostic tests (RDTs) detecting parasite antigens, with PCR available as a reference method for species confirmation.

\section*{Alternative Names}
Malaria Smear; Thick and Thin Blood Films; Blood Smear for Malaria; Malaria RDT; Malaria Antigen Test; Plasmodium PCR; Malarial Parasite Screen; Giemsa Stain for Malaria Parasites

\section*{Principle}
Microscopy (Giemsa-stained thick and thin blood films): Gold standard and WHO-recommended method \cite{who2023malaria}. Thick film: A drop of blood is spread on a slide in a circle ~1 cm in diameter and air-dried (not fixed). Dehemoglobinization occurs during Giemsa staining (the intact RBCs are lysed by the aqueous stain), concentrating parasites 20--40x and allowing detection of low parasitemia (10--50 parasites/$\mu$L by an experienced microscopist). Thick films are used for parasite detection and quantification (parasites/$\mu$L or \% parasitemia). Thin film: A drop of blood is spread in a feathered edge, fixed in methanol, and Giemsa stained. Provides excellent parasite morphology, allowing species identification and determination of \% parasitemia (percentage of parasitized RBCs). Examination of 200 oil-immersion fields at 1,000x magnification is the standard. Rapid diagnostic tests (RDTs): Lateral-flow immunochromatographic tests detecting \textit{Plasmodium} antigens in blood. Three types of targets: (1) HRP-2 (histidine-rich protein 2), specific to \textit{P. falciparum}; (2) pLDH (parasite lactate dehydrogenase), pan-specific and species-specific (\textit{P. falciparum}-specific pLDH and pan-Plasmodium pLDH); (3) Aldolase (pan-specific). Sensitivity approaches microscopy (>95\% for \textit{P. falciparum} at parasitemia >100 parasites/$\mu$L; lower for non-falciparum species). RDTs are particularly important where expert microscopy is unavailable \cite{mathison2017update}. Results in 15--20 minutes. HRP-2 RDTs can remain positive for up to 3--4 weeks after successful treatment (detection of persistent antigen) and should NOT be used to monitor treatment response. \textit{pfhrp2/3} gene deletions (found in some \textit{P. falciparum} strains, particularly in the Horn of Africa and South America) cause false-negative HRP-2-based RDTs. PCR: Real-time PCR targeting \textit{Plasmodium} 18S rRNA gene for genus-level detection with species-specific probes or melting curve analysis. Sensitivity 1--5 parasites/$\mu$L (exceeding expert microscopy). Used in reference laboratories for species confirmation, detection of mixed infections, and when microscopy and RDT are negative but clinical suspicion persists.

\section*{History}
Malaria has been recognized since antiquity (described in Chinese medical texts from 2700 BCE and by Hippocrates). The \textit{Plasmodium} parasite was first visualized in blood by Charles Laveran in 1880 (Nobel Prize 1907). The mosquito transmission cycle was elucidated by Ronald Ross in 1897 (Nobel Prize 1902). The Giemsa stain for blood smears was developed by Gustav Giemsa in 1904 and remains essentially unchanged. Chloroquine, the first synthetic antimalarial, was developed in 1934. Artemisinin, derived from the Chinese herb \textit{Artemisia annua} (sweet wormwood), was discovered by Tu Youyou in the 1970s (Nobel Prize 2015). Artemisinin-based combination therapies (ACTs) became the WHO-recommended first-line treatment for uncomplicated falciparum malaria in the 2000s. Artemisinin resistance (partial, associated with \textit{pfk13} mutations leading to delayed parasite clearance) emerged in the Greater Mekong Subregion in the 2000s and has more recently been reported in Africa (Rwanda, Uganda), threatening global malaria control \cite{ashley2014spread}. RDTs were developed in the 1990s and scaled up in the 2000s, allowing parasitologic diagnosis in resource-limited settings. The WHO's Global Technical Strategy for Malaria (2016--2030) aims for a 90\% reduction in malaria incidence and mortality by 2030.

\section*{Why This Test Is Ordered}
\begin{itemize}
  \item Evaluation of acute febrile illness in a traveler returning from a malaria-endemic area (any fever in a traveler from a malaria-endemic area is malaria until proven otherwise). Testing should be performed STAT
  \item Diagnosis of malaria in a patient with paroxysmal fever, rigors, headache, myalgias, anemia, thrombocytopenia, and splenomegaly, regardless of travel history, though malaria in the absence of travel is exceedingly rare in non-endemic countries unless transfusion- or transplant-associated
  \item Investigation of fever of unknown origin in an immigrant or refugee from a malaria-endemic area, particularly with anemia and thrombocytopenia
  \item Monitoring of parasitemia in patients under treatment for severe falciparum malaria: serial blood smears (or RDTs with microscopy) every 12--24 hours to confirm parasite clearance; parasitemia should decline by >75\% at 48 hours and be negative by day 7 in drug-susceptible infection
  \item Screening of blood donors who have traveled to malaria-endemic regions (deferred per FDA guidance; serologic screening for antibodies plus travel history, not blood smear)
  \item Evaluation of congenital malaria in neonates born to mothers from malaria-endemic areas presenting with fever, irritability, hepatosplenomegaly, anemia, and jaundice in the first 2 months of life
\end{itemize}

\section*{Primary vs Secondary Test}
Microscopy of Giemsa-stained thick and thin blood films is the primary diagnostic test for malaria and should be ordered STAT in suspected cases. For \textit{P. falciparum}, diagnostic delay of even hours can be fatal. RDT is a primary point-of-care test where expert microscopy is not available. PCR is a secondary confirmatory test (reference laboratory) for species identification and when microscopy/RDT are negative but clinical suspicion remains.

\section*{How This Test Is Performed}
Blood films: Venous blood (EDTA/purple-top tube) or fingerstick capillary blood is used to prepare both thick and thin films on glass slides. The thin film is fixed in anhydrous methanol and the thick film is left unfixed. Both are stained with Giemsa (or Wright-Giemsa) at pH 7.2. At least 200 oil-immersion fields (1,000x) on the thick film should be examined before calling the smear negative (in non-immune patients). The level of parasitemia is reported as a percentage of parasitized red blood cells (on thin film) OR parasites per microliter (on thick film, based on counting parasites against white blood cells, assuming 8,000 WBC/$\mu$L). For \textit{P. falciparum}, \% parasitemia >5\% or >100,000 parasites/$\mu$L defines severe malaria. The species should be identified based on thin-film morphology (ring forms, gametocytes, Schizonts, stippling, RBC size, and presence of multiply infected RBCs). RDT: A drop of fingerstick or venous blood is applied to the test strip following the manufacturer's instructions; results in 15--20 minutes. PCR: EDTA whole blood (200 $\mu$L); should be sent to a reference laboratory (CDC, state public health laboratory).

\section*{What Preparation Is Needed}
No special patient preparation. The timing of blood collection relative to fever paroxysms: Parasitemia is typically highest during or shortly after a fever spike. However, blood films should be obtained immediately when malaria is suspected, NOT delayed to coincide with a fever spike. If the initial smear is negative and clinical suspicion is high, blood films should be repeated every 12--24 hours for a total of 3 sets (one set = thick + thin films). This is because parasitemia in non-immune patients can be very low initially or may be sequestered (cytoadherence of \textit{P. falciparum}-infected RBCs in the microvasculature).

\section*{Contraindications and Precautions}
A single negative blood film does NOT exclude malaria, particularly in non-immune patients (who may become symptomatic at very low parasitemias of <50 parasites/$\mu$L, at the limit of detection for thick-film microscopy). If clinical suspicion persists, repeat smears every 12--24 hours for at least 3 sets. RDTs have reduced sensitivity for non-falciparum species (especially \textit{P. ovale} and \textit{P. malariae}) and for low parasitemias (<100 parasites/$\mu$L). HRP-2-based RDTs will remain positive for 2--4 weeks after successful treatment (persistent antigenemia); a positive RDT after treatment does NOT indicate treatment failure. \textit{P. knowlesi} appears morphologically similar to \textit{P. malariae} on microscopy (band-form trophozoites) but can cause severe, rapidly progressive falciparum-like disease and must be confirmed by PCR. In \textit{P. falciparum} malaria, only ring forms (trophozoites) and gametocytes are seen in peripheral blood (mature trophozoites and schizonts sequester in the microvasculature). The presence of mature forms in peripheral blood suggests \textit{P. falciparum} with very high parasitemia or misidentification of a non-falciparum species. Immune individuals (residents of endemic areas) may have parasitemia without clinical symptoms; their smears should be interpreted with caution, and a positive smear does not necessarily mean acute malaria is the cause of their current illness.

\section*{Risks}
Venipuncture carries standard risks. There is no risk from the test itself. The major clinical risk is failure to diagnose \textit{P. falciparum} malaria promptly. \textit{P. falciparum} malaria can progress from uncomplicated to severe disease (coma, renal failure, ARDS, shock) within hours, and every hour of delay in effective antimalarial therapy increases mortality.

\section*{What Happens After the Test}
Blood smear results should be available STAT (within 1--2 hours). RDT results in 15--20 minutes. PCR results in 24--72 hours at a reference laboratory. If the smear or RDT is positive: (1) Identify species and determine \% parasitemia. (2) For \textit{P. falciparum} (or species not yet determined in a severely ill patient): initiate urgent treatment. Uncomplicated \textit{P. falciparum} (acquired in chloroquine-sensitive area---Haiti, Central America west of the Panama Canal, parts of the Middle East): artemether-lumefantrine or atovaquone-proguanil. Uncomplicated \textit{P. falciparum} (chloroquine-resistant area): ACT (artemether-lumefantrine preferred in the US) or atovaquone-proguanil. Severe \textit{P. falciparum} (any parasitemia $\geq$5\%, or criteria for severe malaria: cerebral malaria, ARDS, renal failure, shock, DIC, acidosis, hypoglycemia, jaundice): IV artesunate (obtained from CDC under an IND protocol, as it is not FDA-approved or commercially available in the US) is first-line and superior to IV quinidine/quinine. If IV artesunate cannot be obtained immediately, IV quinidine (plus doxycycline/tetracycline/clindamycin) should be initiated while awaiting artesunate. (3) Admit all patients with \textit{P. falciparum} malaria for at least 24--48 hours to monitor for clinical deterioration. For \textit{P. vivax} or \textit{P. ovale}: chloroquine (or ACT) for the blood-stage infection PLUS primaquine or tafenoquine for radical cure of hypnozoites (must test G6PD before primaquine/tafenoquine due to risk of severe hemolysis in G6PD deficiency) \cite{cdc2023treatment}.

\section*{Understanding Results}

\subsection*{Below Normal}
\begin{itemize}
  \item \textbf{Blood film negative (no malaria parasites seen):} No \textit{Plasmodium} parasites detected in the thick film after examining 200--300 oil-immersion fields. A single negative smear does NOT exclude malaria. Sensitivity depends on parasitemia level: at 100 parasites/$\mu$L, sensitivity approaches 100\% with expert microscopy; at 10--50 parasites/$\mu$L, sensitivity is ~90\%. Repeat smears every 12--24 hours for at least 3 sets if clinical suspicion persists
  \item \textbf{RDT negative:} No HRP-2, pLDH, or aldolase detected. Does not exclude non-falciparum malaria (low sensitivity for \textit{P. ovale, P. malariae}) or low-level \textit{P. falciparum} parasitemia. False-negative HRP-2 RDT can occur with \textit{pfhrp2/3}-deleted \textit{P. falciparum} (prevalent in the Amazon and Horn of Africa). PCR or repeat microscopy is indicated if suspicion remains
  \item \textbf{PCR negative:} \textit{Plasmodium} DNA not detected. Essentially excludes malaria
\end{itemize}

\subsection*{Above Normal}
\begin{itemize}
  \item \textbf{Blood film positive for \textit{Plasmodium}:} Confirms malaria. Species determination is critical:
    \begin{itemize}
      \item \textbf{\textit{P. falciparum}:} Ring forms (small, delicate rings, often multiple per RBC, accole/appliqué forms at the RBC margin), banana-shaped gametocytes (pathognomonic), no trophozoites or schizonts in peripheral blood (unless very high parasitemia). High \% parasitemia ($\geq$5\%) defines severe malaria. Treatment: ACT or atovaquone-proguanil for uncomplicated; IV artesunate for severe
      \item \textbf{\textit{P. vivax}:} Enlarged parasitized RBCs (1.5--2x normal), Schüffner stippling (fine red/pink dots), ameboid ring forms, all stages in peripheral blood. Associated with hypnozoite relapses. Requires primaquine/tafenoquine for radical cure (check G6PD first)
      \item \textbf{\textit{P. ovale}:} Similar to \textit{P. vivax} (enlarged RBCs, Schüffner stippling), but more compact trophozoites. Also has hypnozoites. Requires primaquine/tafenoquine for radical cure after G6PD testing
      \item \textbf{\textit{P. malariae}:} Normal-sized or small RBCs, band-form trophozoites (pathognomonic), rosette schizonts (8--10 merozoites). Chronic infection can last decades. No hypnozoites. Chloroquine is first-line
      \item \textbf{\textit{P. knowlesi}:} Morphologically mimics \textit{P. malariae} (band forms) but can cause severe falciparum-like disease. Confirmed by PCR. Transmitted from macaques in SE Asia
    \end{itemize}
  \item \textbf{RDT positive:} Confirms malaria at the bedside. HRP-2 positive = \textit{P. falciparum}. Pan-pLDH or aldolase positive = non-falciparum \textit{Plasmodium}. Mixed infection (HRP-2+ plus pan-pLDH+) suggests \textit{P. falciparum} and a non-falciparum co-infection
  \item \textbf{Severe malaria criteria (WHO):} Any one of: impaired consciousness (GCS <11 or Blantyre coma score <3 in children), prostration, multiple convulsions, acidotic breathing, hypoglycemia (<40 mg/dL), severe anemia (Hb <7 g/dL in adults, <5 g/dL in children), renal impairment (creatinine >3.0 mg/dL or urine output <400 mL/24h), jaundice (bilirubin >3.0 mg/dL), pulmonary edema/ARDS, significant bleeding/DIC, shock (systolic BP <80 mmHg in adults), hyperparasitemia ($\geq$5\%)
\end{itemize}

\section*{Normal Results}
\textbf{No malaria parasites seen on Giemsa-stained thick and thin blood films (examined 200--300 oil-immersion fields).} RDT: \textbf{Negative.} PCR: \textbf{\textit{Plasmodium} DNA not detected.}

\section*{Follow-up Tests}
\begin{itemize}
  \item \textbf{Repeat blood smears every 12--24 hours (up to 3 sets):} If initial smears are negative but clinical suspicion remains high. A single negative smear does NOT exclude malaria
  \item \textbf{Malaria PCR (Plasmodium 18S rRNA gene):} For species confirmation when microscopy is equivocal, for distinguishing \textit{P. knowlesi} from \textit{P. malariae}, or when smears are negative but suspicion persists. CDC's Malaria Branch provides PCR for reference diagnostic confirmation (assay detects all five human-infecting species)
  \item \textbf{CBC with platelet count:} Anemia (hemoglobin typically 7--10 g/dL, worsening with parasitemia), thrombocytopenia (platelets often <100,000/$\mu$L, a sensitive but nonspecific finding in malaria), and leukopenia are characteristic. Thrombocytopenia is present in 60--80\% of malaria cases regardless of species
  \item \textbf{Comprehensive metabolic panel:} For evidence of severe malaria: renal function (creatinine, BUN), liver function (bilirubin, ALT/AST --- typically elevated), glucose (hypoglycemia is a common and dangerous complication, particularly with IV quinine/quinidine), and lactate (acidosis)
  \item \textbf{G6PD (glucose-6-phosphate dehydrogenase) activity:} Before administering primaquine or tafenoquine for radical cure of \textit{P. vivax} or \textit{P. ovale}, G6PD deficiency must be excluded to prevent severe iatrogenic hemolysis. Quantitative G6PD enzyme activity is preferred
  \item \textbf{Blood cultures:} To exclude concurrent bacterial sepsis, which can mimic or complicate severe malaria (particularly in children in endemic areas where bacteremia co-infection is common)
\end{itemize}

\section*{References}
\bibliographystyle{unsrtnat}
\bibliography{references}

\clearpage
\section*{Regulatory Status and Authorization Date}
Regulatory status applies to the specific assay, instrument, manufacturer, and intended use rather than to the test name alone. No single approval date can be assigned to this test category. For a commercial assay, verify the current product in the relevant regulator's database: the U.S. Food and Drug Administration (FDA) clearance, approval, or authorization; India's Central Drugs Standard Control Organisation (CDSCO) licence or permission; the European Union In Vitro Diagnostic Regulation (EU IVDR) CE certificate issued through the applicable conformity-assessment route; or the United Kingdom Medicines and Healthcare products Regulatory Agency (MHRA) registration and UKCA/CE status. Laboratory-developed tests may not have an individual FDA, CDSCO, EU IVDR, or MHRA approval date.
\clearpage
